\section{Limitations}
\label{sec:limitations}

\paragraph{No real-device measurements.} This paper reports an emulator reproduction, a prototype measurement on the same emulator, and a clearly labeled simulation; it contains no device latency, energy, or hardware-keystore numbers. On an emulator or Cuttlefish image the keystore is software-only, so we report \texttt{security level = software} and make no hardware-protection claim. Device measurements are future work.

\paragraph{Prototype, not the \texttt{system\_server} design.} The B3 measurement comes from a grant-based prototype: a notification listener, an in-process vault, and an accessibility overlay. It therefore covers the notification channel and the screenshot channel only. The accessibility tree still exports the underlying text to an agent that reads it directly; the clipboard is not mediated (our agent did try long-press, copy, and paste); and a production design would redact node text in the framework instead of covering pixels. The reveal is a plain button rather than a user-presence gate, and \texttt{FLAG\_SECURE} prevents capture, not accessibility-tree reads. With 30 trials per condition, one emulator image, and one model, the $29/30$ versus $0/30$ comparison is a directional signal, not a device result.

\paragraph{Binding is not honest use.} Destination binding constrains \emph{where} a value may be resolved, not how a trusted app then uses it. An injected instruction can still cause an agent to perform a legitimate-but-unwanted action inside an allowed destination; we scope this out and quantify it as a residual risk in the evaluation plan.

\paragraph{WebView fail-closed.} Binding a field to a WebView document origin requires node-to-frame-to-origin attestation, which we do not solve here; web fields are rejected rather than bound. This trades some utility for honesty.

\paragraph{Screenshot-only agents.} An agent that can only act on pixels without naming a stable target node cannot be given a bound field; such agents need a trusted target-selection adapter or user confirmation and are out of scope.

\paragraph{Simulation is not evidence.} The reference model demonstrates design trends under synthetic canaries. It does not establish the absence of exposure on uninstrumented channels and is not a substitute for the device evaluation.

\paragraph{Conditional guarantees.} The results of Section~\ref{sec:model} hold only under the reference-monitor assumptions and characterize the model, not an implementation. They are not measurements, are not claimed tight or optimal, and can fail silently if any assumption is violated---for example, if an observation channel outside the modeled set is not redacted, or if a fallback copies the value to the clipboard or asks the agent to type it, which invalidates the guarantee for that operation.

\paragraph{Composition gaps.} A deployment combines several sources, an agent framework, input components, OEM services, and applications. Non-disclosure composes only if every component that receives the value is either inside the trust boundary or outside the agent view, and destination confinement composes only if every resolution path uses the same target identity. We do not yet have a mechanism that proves an arbitrary third-party keyboard, accessibility overlay, or OEM modification is in the modeled set, so the prototype reports which channels it instruments and treats the rest as assumptions rather than claims.
